\section{Contact formulation and discretization}\label{sec:form}

The transition map of Section~\ref{sec:tmap} acts as a geometric oracle: queried at any point, it
reports two numbers that a contact solver needs and nothing more. The first is the gap $\gn$, the
signed distance from the point to the obstacle surface. The second is the outward unit normal $\nbf$,
the direction along which contact pushes. The atlas reads this pair from the boundary-to-boundary
chart map $\tmap=\map{B}^{-1}\circ\map{A}$, which composes the inverse of one body's boundary chart
with the other's. A neural signed-distance function (a level set) supplies the same pair as a
comparison baseline. This section places that oracle inside a standard contact solver and keeps the
two sources interchangeable.

Everything the solver does with $(\gn,\nbf)$ is classical, and we want the reader to see that nothing
in it depends on where the pair came from. The constraint enforcement (penalty and augmented
Lagrangian), the friction law, the consistent Newton tangent, and the spatial discretization are all
textbook constructions. We route them through a single interface that returns $(\gn,\nbf)$, so that
force assembly, broad phase, tangent, and well-posedness are byte-for-byte identical whether the
geometry comes from the atlas or from the level set. The accuracy gap reported in
Sections~\ref{sec:verif}--\ref{sec:cv7} is therefore charged to the geometry alone, not to the
constraint solver.

We fix the sign convention once. The gap $\gn(\xbf)$ is the signed distance from a candidate contact
point $\xbf$ to the obstacle surface: positive when the bodies are apart, negative when they overlap,
so $\gn<0$ flags penetration. The normal $\nbf$ points outward from the obstacle, which makes the
contact force $-t_N\nbf$ push a penetrating point back out, where $t_N\ge0$ is the contact pressure.
We write $\pos{\cdot}$ for the Macaulay bracket $\pos{a}=\max(0,a)$, which returns a quantity only
when it is positive, and zero otherwise.

\subsection{Signorini conditions and the contact potential}\label{sec:form-signorini}

% code: solvers/contact/augmented_lagrangian.py
Contact is governed by three plain physical rules: bodies may not interpenetrate, contact only
pushes and never pulls, and pressure builds up only where the surfaces actually touch. We state them
for two deformable bodies $\Omega^A,\Omega^B\subset\R^3$ that undergo displacement fields
$\ubf^A,\ubf^B$ and may meet along a candidate contact surface $\Gamma_c$ in the deformed
configuration. The three rules become the Signorini--Hertz--Moreau conditions, which are the
Karush--Kuhn--Tucker (KKT) optimality conditions of the constrained problem,
\begin{equation}\label{eq:signorini}
\gn\ \ge\ 0,\qquad t_N\ \ge\ 0,\qquad \gn\, t_N\ =\ 0\qquad\text{on }\Gamma_c .
\end{equation}
Reading them in order: the first inequality forbids penetration; the second forbids tensile
(adhesive) tractions, so the surface can push but not glue; and the complementarity $\gn t_N=0$ ties
the two together, allowing pressure only where the gap has closed to zero. These unilateral
conditions of \citet{johnson1985contact} and \citet{wriggers2006computational} are what make
equilibrium a variational inequality rather than a smooth minimization
\citep{kikuchi1988contact,laursen2002computational}: the active contact set is itself an unknown.

We now write equilibrium in weak (virtual-work) form so that the contact pressure enters as a force
on the boundary. Let
$a^\alpha(\ubf^\alpha,\delta\ubf^\alpha)=\int_{\Omega^\alpha}\bm{P}(\bm{F}^\alpha):\grad_{\xbf}\delta\ubf^\alpha\,dV$
be the internal virtual work of body $\alpha\in\{A,B\}$, in which $\bm{P}$ is the first
Piola--Kirchhoff stress, $\bm{F}^\alpha=\bm{I}+\grad\ubf^\alpha$ is the deformation gradient, and
$\delta\ubf^\alpha$ is an admissible virtual displacement; let $L^\alpha$ collect the external
virtual work of the body and surface loads. We need one more ingredient: how the gap responds to a
virtual displacement. Because the gap is a distance measured along the normal, its first variation is
just the relative virtual motion projected onto $\nbf$,
\begin{equation}\label{eq:gap-variation}
\delta\gn\ =\ \nbf\cdot(\delta\ubf^B-\delta\ubf^A),
\end{equation}
which uses only the contact normal and is exact to first order no matter how the surface is
represented. This is the key point for the atlas-versus-level-set comparison: the variation depends on
the geometry through $\nbf$ alone. Constrained equilibrium then reads: find $(\ubf^A,\ubf^B)$ and a
pressure multiplier $t_N\ge0$ such that
\begin{equation}\label{eq:weak-contact}
\sum_{\alpha\in\{A,B\}}\!\Big[a^\alpha(\ubf^\alpha,\delta\ubf^\alpha)-L^\alpha(\delta\ubf^\alpha)\Big]
\ +\ \int_{\Gamma_c} t_N\,\delta\gn\,dA\ =\ 0,
\end{equation}
holds together with the Signorini conditions~\eqref{eq:signorini}. The new term is the contact
integral: the pressure $t_N$ does virtual work against the gap variation along $\Gamma_c$. What
remains is to supply $t_N$, and the next two paragraphs give the two standard ways to do so.

\paragraph{Penalty regularization: a stiff one-sided spring.}
% code: solvers/contact/penalty.py::compute_contact_force
The penalty method replaces the exact pressure by a stiff spring that switches on only when the
constraint is violated. Picture a spring that does nothing while the bodies are apart and pushes back
ever harder the deeper they overlap. With penalty stiffness $\varepsilon_n>0$ the pressure is
\begin{equation}\label{eq:penalty-pressure}
t_N\ =\ \varepsilon_n\,\pos{-\gn},
\end{equation}
which is zero whenever $\gn\ge0$ and rises linearly with the penetration depth $-\gn$ once the bodies
overlap. The spring stores energy as it compresses, and that stored energy is the smooth contact
potential,
\begin{equation}\label{eq:penalty-potential}
\Pi_c(\gn)\ =\ \tfrac{1}{2}\varepsilon_n\,\pos{-\gn}^2,
\qquad
t_N\ =\ \frac{\partial \Pi_c}{\partial(-\gn)}\ =\ \varepsilon_n\,\pos{-\gn} .
\end{equation}
Differentiating the potential recovers the pressure~\eqref{eq:penalty-pressure}, which tells us the
normal contact force is conservative: it is the gradient of a scalar energy. Because that energy is
non-negative, adding it to the total potential keeps the problem coercive, the property used in
Section~\ref{sec:form-wellposed}. The force carried by one quadrature point (or, in the meshfree
setting, one particle) in physical space then follows from the pressure,
\begin{equation}\label{eq:penalty-force}
\fbf^{\,\mathrm{N}}\ =\ t_N\,A\,\nbf\ =\ \varepsilon_n\,\pos{-\gn}\,A\,\nbf ,
\end{equation}
where $A$ is the tributary area carried by that point (or the particle volume $V_p$ in the MPM
coupling of Section~\ref{sec:form-mpm}). The force points along the normal and scales with both the
penetration and the stiffness. Equation~\eqref{eq:penalty-force} is exactly what
\code{compute\_contact\_force} in \code{solvers/contact/penalty.py} does: clamp $-\gn$ to its positive
part, scale by $\varepsilon_n A$, and project onto the obstacle normal. We size the spring to the
material it restrains, $\varepsilon_n=\beta E/h$, with $\beta\in[1,20]$, $E$ the Young modulus, and
$h$ the local mesh size, so the contact stiffness matches the bulk stiffness. The leftover penetration
then scales as $\delta\sim\sigma_n/\varepsilon_n$ and shrinks as $\varepsilon_n$ grows.

\paragraph{Augmented Lagrangian: a spring whose rest force is remembered.}
% code: solvers/contact/augmented_lagrangian.py
The penalty spring forces a compromise: a stiffer $\varepsilon_n$ shrinks the residual penetration but
worsens the conditioning of the system, and in explicit dynamics it shortens the stable time step. The
augmented-Lagrangian method of \citet{alart1991mixed} and \citet{simo1992augmented} escapes the
compromise by giving the spring a memory. Alongside the penalty it carries a persistent multiplier
$\lambda\ge0$ that remembers the pressure accumulated so far, so the spring no longer has to supply the
full pressure from penetration alone. The augmented pressure and its Uzawa update are
\begin{equation}\label{eq:al-pressure}
t_N\ =\ p_{\mathrm{aug}}(\gn,\lambda)\ =\ \pos{\lambda-\varepsilon_n\,\gn},
\qquad
\lambda^{k+1}\ =\ \pos{\lambda^{k}-\varepsilon_n\,\gn(\ubf^{k+1})} .
\end{equation}
Reading the update: with $\gn<0$ on overlap, the term $-\varepsilon_n\gn$ is positive, so both the
pressure and the remembered multiplier grow while penetration persists; once the gap closes
($\gn\ge0$) the Macaulay clamp drives both back to zero. The payoff is that at convergence the memory
$\lambda$ settles on the true contact pressure, $\lambda\to t_N^\star$, and the penetration vanishes
without sending $\varepsilon_n\to\infty$ \citep{laursen2002computational}. This is
\code{AugmentedLagrangianContact} in \code{solvers/contact/augmented\_lagrangian.py}; the solver stores
$\lambda$ per contact point and reuses it across iterations, so the caller must pass the gap with
stable indexing (the implementation resets $\lambda$ and warns if the array shape changes between
calls). The plain penalty~\eqref{eq:penalty-pressure} is the special case with no memory,
$\lambda\equiv0$.

\subsection{Regularized Coulomb friction}\label{sec:form-friction}

% code: solvers/contact/friction.py
Once contact carries a normal pressure, friction resists tangential sliding. The Coulomb law caps the
tangential traction $\tbf_T$ by the friction cone $\norm{\tbf_T}\le\mu\,t_N$: inside the cone the
surfaces stick ($\vbf_T=0$), and on its boundary they slip with
$\tbf_T=-\mu t_N\,\vbf_T/\norm{\vbf_T}$, where $\mu$ is the friction coefficient and $\vbf_T$ the
tangential slip velocity \citep{wriggers2006computational}. This law has a kink at zero sliding speed
that a Newton solver cannot differentiate, so we smooth it. We replace the sharp law by the
bracket form
\begin{equation}\label{eq:friction-reg}
\tbf_T\ =\ -\,\mu\,\norm{\fbf^{\,\mathrm{N}}}\,
\frac{\vbf_T}{\sqrt{\norm{\vbf_T}^2+\varepsilon_T^2}},
\end{equation}
in which the tangential velocity is the full velocity with its normal component removed,
\begin{equation}\label{eq:vt-projection}
\vbf_T\ =\ \big(\bm{I}-\nbf\otimes\nbf\big)\,\vbf\ =\ \vbf-(\vbf\cdot\nbf)\,\nbf .
\end{equation}
The parameter $\varepsilon_T$ is a small velocity scale that controls the rounding of the kink. When
sliding is fast, $\norm{\vbf_T}\gg\varepsilon_T$, the traction approaches the cone bound
$\mu\norm{\fbf^{\,\mathrm{N}}}$ and the law behaves like sharp Coulomb slip; when sliding is slow,
$\norm{\vbf_T}\ll\varepsilon_T$, the traction grows linearly with $\vbf_T$ and passes smoothly through
zero, so stick is approximated by a stiff tangential spring. Like the penalty, this surrogate comes
from a potential: the pseudo-potential
$\Psi_T(\vbf_T)=\mu\,t_N\sqrt{\norm{\vbf_T}^2+\varepsilon_T^2}$ yields
$\tbf_T=-\partial\Psi_T/\partial\vbf_T$, so the regularized problem stays variational. In code,
\code{compute\_friction\_force} in \code{solvers/contact/friction.py} takes the scalar normal-force
magnitude $\norm{\fbf^{\,\mathrm{N}}}$ as input, so it composes unchanged with either normal model: the
penalty pressure ($\norm{\fbf^{\,\mathrm{N}}}=\varepsilon_n\pos{-\gn}A$) or the augmented-Lagrangian
pressure ($\norm{\fbf^{\,\mathrm{N}}}=p_{\mathrm{aug}}A$). The total contact force is then the sum of
normal and tangential parts, $\fbf^{\,\mathrm{contact}}=\fbf^{\,\mathrm{N}}+\tbf_T A$. By construction
the friction force opposes the slip, $\tbf_T\cdot\vbf_T\le0$, which the verification suite confirms
node by node.

The quasi-static rough-joint driver of Section~\ref{sec:cv7} has no sliding velocity to regularize, so
there we drive friction differently. We drag each active node along a prescribed full-slip direction
$\tbf_{\mathrm{proj}}=\widehat{\hat{\tbf}-(\hat{\tbf}\cdot\nbf)\nbf}$, formed by projecting the loading
direction $\hat{\tbf}$ onto the tangent plane and renormalizing, and set
$\tbf_T=\mu\,t_N\,\tbf_{\mathrm{proj}}$. The velocity surrogate~\eqref{eq:friction-reg} is instead the
path used in the explicit MPM cross-check. Both sit at the cone bound and are plain Coulomb, and
neither imposes a dilation angle. The dilatant, slope-dependent response seen in
Section~\ref{sec:cv7} is therefore not put in by hand. It emerges from the resolved asperity geometry
through the normal $\nbf$, rather than from a constitutive dilation law of the
Patton~\citep{patton1966multiple} or Barton~\citep{barton1977shear,barton1990review} type.

\subsection{Consistent contact tangent}\label{sec:form-tangent}

% code: benchmarks/contact/cv_numerical/rock_joint_decoder_shear.py::_contact_tangent
A Newton solver needs the derivative of the contact force with respect to the displacement; supplying
the right derivative is what makes Newton converge quadratically. Writing the residual as
$\bm{R}(\ubf)=\bm{K}_{\mathrm{elastic}}\ubf-\fbf_{\mathrm{contact}}(\ubf)$, the contact part of the
tangent at an active node ($\gn<0$) gathers three contributions,
\begin{equation}\label{eq:contact-tangent}
\bm{K}_c\ =\ \frac{\partial \fbf_{\mathrm{contact}}}{\partial \ubf}
\ =\ \underbrace{\varepsilon_n A\,\big(\nbf\otimes\nbf\big)}_{\text{normal penalty stiffness}}
\ +\ \underbrace{\varepsilon_n A\,\mu\,\big(\tbf_{\mathrm{proj}}\otimes\nbf\big)}_{\text{friction--normal coupling}}
\ +\ \underbrace{t_N A\,\frac{\partial\nbf}{\partial\ubf}}_{\text{curvature / shape operator}} .
\end{equation}
Each term answers a distinct question. The first is the normal stiffness: it is the penalty stiffness
$\varepsilon_n A$ projected onto the active normal through $\nbf\otimes\nbf$, and it tells Newton how
the normal force stiffens as the node sinks in. Because $\nbf$ tilts with the asperity slope,
$\nbf\otimes\nbf$ is a full rank-one $3\times3$ block, not the single vertical stiffness a
flat-interface model would assume, so a tilted asperity already couples the vertical and horizontal
directions. The second term is the friction--normal coupling: it tells Newton how the frictional drag
changes when the normal pressure changes. It is non-symmetric, which is the structural signature of a
non-associative friction law \citep{wriggers2006computational}. The third term is the curvature (or
geometric) stiffness: it accounts for the normal direction itself rotating as the surface deforms,
which matters wherever the surface bends. That rotation is the shape operator (the Weingarten map) of
the contact surface,
\begin{equation}\label{eq:shape-operator}
\frac{\partial n_i}{\partial x_j}
\ =\ \frac{1}{\norm{\grad\phi}}\,\big(\delta_{ik}-n_i n_k\big)\,
\frac{\partial^2\phi}{\partial x_k\partial x_j},
\end{equation}
which the surface description supplies through its second derivative: either second-order autograd on
the neural SDF, or the second derivative $h_\theta''$ of the Fourier-feature height chart. In
practice, the rough-joint driver assembles the first two terms (this is
\code{\_contact\_tangent}). On the smooth single-asperity benchmarks the curvature term is small, so
keeping only the rank-one normal and friction--normal blocks, with an Armijo line search, still drives
the frictional Newton residual down to $\sim10^{-9}$ on the one-block shear, a fully converged solve.
One configuration resists this convergence. On the two-deformable, moving-master node-to-surface
problem the friction residual settles at $1.5$--$6\%$ and does not shrink under mesh refinement,
because the contact point itself moves and makes the active set non-smooth. The remedy there is a
semismooth-Newton, augmented, or mortar treatment
\citep{alart1991mixed,puso2004mortar,chouly2013nitsche}, not a finer mesh.

\subsection{Discretization: chart-FEM pushforward and physical-space MPM}\label{sec:form-mpm}

We enforce the contact constraint on two discretizations: a chart-based finite element method and a
chart-based material point method. Both share the same force expressions and the same $(\gn,\nbf)$
oracle. They differ only in how they carry the geometry of the chart, and the contrast between the two
is the point of this subsection.

\paragraph{Chart-FEM pushforward.}
% code: solvers/fem/chart_vector_fem.py
The chart-FEM solves the elasticity problem on a simple reference cube and then maps it onto the curved
physical body. Concretely, on each chart the elastostatic problem lives on the reference parametric
cube $U_i=[-r,r]^3$, and the chart decoder $\map{i}$ carries it to physical space through its Jacobian
$\bm{J}_i=\partial\map{i}/\partial\bm{\xi}$, the matrix of derivatives of the physical coordinate with
respect to the reference coordinate $\bm{\xi}$. Gradients computed on the cube must be converted to
physical gradients by the chain rule, which means multiplying by $\bm{J}_i^{-1}$,
\begin{equation}\label{eq:fem-pushforward}
\grad_{\xbf}\ubf\ =\ \grad_{\bm{\xi}}\ubf\,\bm{J}_i^{-1},
\qquad
\bm{F}\ =\ \bm{I}+\grad_{\bm{\xi}}\ubf\,\bm{J}_i^{-1},
\end{equation}
so $\grad_{\xbf}\ubf$ is the physical displacement gradient and $\bm{F}$ the physical deformation
gradient. The same chain rule and the change-of-variables volume factor $\abs{\det\bm{J}_i}$ then turn
the cube integral into the physical stiffness,
\begin{equation}\label{eq:fem-assembly}
\bm{K}_{ab}\ =\ \int_{U_i}\big(\grad_{\bm{\xi}}N_a\,\bm{J}_i^{-1}\big)^{\!\top}
:\mathbb{C}:\big(\grad_{\bm{\xi}}N_b\,\bm{J}_i^{-1}\big)\,\abs{\det\bm{J}_i}\,d\bm{\xi},
\end{equation}
where $N_a,N_b$ are the shape functions and $\mathbb{C}$ is the elasticity tensor. The one subtlety is
that a neural chart can be steep, which makes $\bm{J}_i$ nearly singular and $\bm{J}_i^{-1}$ large. We
therefore form $\bm{J}_i^{-1}$ and $\abs{\det\bm{J}_i}$ through the SVD-clamped operator
\code{stabilized\_jacobian\_ops} (\code{solvers/fem/chart\_vector\_fem.py}), which floors the singular
values so a poorly conditioned chart cannot blow up the assembly, while $\det\bm{J}_i>0$ (no fold-over)
is monitored throughout. On the boundary-fitted rough geometry this pushforward passes a patch test and
recovers the manufactured-solution rate $O(h^2)$ (measured exponents $2.07$ and $1.83$), which
confirms that the discretization error is controlled independently of the contact treatment. The
penalty force $\fbf^{\,\mathrm{N}}=\varepsilon_n\pos{-\gn}A\,\nbf$ enters the chart residual
work-conjugately, with $A$ the tributary boundary-face area measured in physical space.

\paragraph{MPM coupling in physical space.}
% code: docs/contact_theory_manual.md section 1.2
The material point method needs no such pushforward for the contact force, and this is the point to
isolate. The chart-based MPM stores particle velocity, gravity, and stress directly in physical
space, so the contact force is added through the same particle-to-grid (P2G) channel as gravity,
with \emph{no Jacobian pull-back},
\begin{equation}\label{eq:mpm-p2g}
\fbf_I^{\mathrm{ext}}\ \mathrel{+}=\ \sum_p \fbf_p^{\,\mathrm{contact}}\,N_I(\bm{\xi}_p),
\qquad
\fbf_p^{\,\mathrm{contact}}\ =\ \varepsilon_n\,\pos{-\gn(\xbf_p)}\,V_p\,\nbf_p\ +\ \tbf_T\,V_p ,
\end{equation}
where $N_I$ is the grid shape function and $\xbf_p=\map{i}(\bm{\xi}_p)+\ubf_p$ is the particle's
physical position. The volume factor $V_p$ converts the penalty pressure into a force (pressure times a
volume), so $\fbf_p^{\,\mathrm{contact}}$ scatters to the grid like any other force: without the mass
weight that gravity carries, and without the $\bm{J}_i^{-\top}$ that a chart-local force would need.
Put plainly, forces and grid velocities already live in the same physical space, so the contact force
is added verbatim, and that is the central bookkeeping point of the MPM coupling. The penalty spring
also adds its own stability bound to the explicit time step,
$\Delta t_{\mathrm{contact}}\le 0.5\sqrt{m_{\min}/\varepsilon_n}$, so the effective step is the smaller
of the elastic CFL limit and this contact limit, $\min(\Delta t_{\mathrm{CFL}},\Delta
t_{\mathrm{contact}})$. The MPM path reproduces the Coulomb floor, returning an apparent friction
$\muapp\approx0.34$ against an imposed $\mu=0.4$, but it does not reproduce the dilatancy, because a
mated rough-on-rough block is unstable in explicit penalty dynamics. We therefore use the MPM as a
dynamic cross-check rather than as the primary rough-joint solver, and we mark that limit plainly.

\subsection{Level-set-free detector dispatch}\label{sec:form-dispatch}

% code: solvers/contact/contact_manager.py::body_gap_normal
The geometry enters the solver at exactly one place, by design. A single function
\code{body\_gap\_normal} takes a batch of physical-space query points and returns the pair
$(\gn,\nbf)$, choosing its source according to the body's declared detector,
\begin{equation}\label{eq:detector-dispatch}
(\gn,\nbf)\ =\
\begin{cases}
\big(\phi_A(\xbf),\ \grad\phi_A/\norm{\grad\phi_A}\big), & \text{detector} = \texttt{sdf},\\[4pt]
\texttt{evaluate\_gap\_chart}(\xbf,\ \text{chart}), & \text{detector} = \texttt{chart}.
\end{cases}
\end{equation}
Both branches honor the identical contract: $\gn$ has shape $(N,)$ with $\gn<0$ meaning penetration,
and $\nbf$ has shape $(N,3)$ and is a unit outward normal. The \texttt{sdf} branch evaluates the
neural level set $\phi_A$ and its autograd gradient; the \texttt{chart} branch evaluates the
transition-map detector of Section~\ref{sec:tmap}, which is the radial chart for star-shaped bodies, or
the height chart $\gn=z-h_\theta(x)$ with normal $(-h_\theta',1)/\sqrt{1+h_\theta'^2}$ for
single-valued joints. Everything downstream reads only $(\gn,\nbf)$: the broad-phase seed-to-seed cull,
the narrow-phase active-set selection $\{\gn<0\}$, the penalty and augmented-Lagrangian forces, the
friction law, the consistent tangent~\eqref{eq:contact-tangent}, and both discretizations. None of it
can tell which branch produced the pair. Swapping the atlas for the level set therefore touches one
branch of~\eqref{eq:detector-dispatch} and nothing else, and that is what makes the comparison
controlled: the same solver, the same loading, the same tolerances, with only the geometric oracle
changed.

% code: solvers/contact/contact_manager.py::body_gap_normal
Figure~\ref{fig:tm-levelset-vs-chart} puts the two oracles side by side on a rough profile. The atlas
detector resolves the asperity slopes; the level set, built from an ambient field whose spectral bias
favors low frequencies \citep{rahaman2019spectral}, smooths them away. The numbers follow the
picture. On the real Inada-granite joint, the chart detector returns a contact gap with $4.2\%$ RMS
error against the analytic surface, while the ambient SDF returns $44.7\%$. That is a factor of $10.5$
on the quantity that drives contact, and the chart reaches it at only $1.47\times$ the per-query cost
of the SDF and with no ambient field to store. That accuracy difference propagates. As
Section~\ref{sec:cv7} quantifies, the level set under-predicts the emergent dilatancy by $98\%$ and the
frictional strength by $35\%$, purely because it feeds the same solver a smoothed normal.

\subsection{Well-posedness and the penalty estimate}\label{sec:form-wellposed}

% code: contact_atlas/03_mathematical_theory.md
We close with the standard existence and stability statements. They carry over unchanged here, because
the contact treatment is classical and the chart changes only how the gap is evaluated, not the
analysis around it. Assume that each body's stored-energy density $W(\bm{F})$ is polyconvex with the
usual coercivity and growth bounds (the Neo-Hookean and small-strain models used here qualify), that
the admissible set $K=\{(\ubf^A,\ubf^B): \gn\ge0\text{ on }\Gamma_c\}$ is non-empty, that the body
loads lie in
$L^2(\Omega)$, and that the Dirichlet data lie in $H^{1/2}(\Gamma_D)$. The natural object to minimize is the
total energy, the stored elastic energy minus the work of the loads,
\begin{equation}\label{eq:total-energy}
E(\ubf^A,\ubf^B)\ =\ \sum_{\alpha\in\{A,B\}}\int_{\Omega^\alpha}W(\bm{I}+\grad\ubf^\alpha)\,dV
\ -\ \sum_{\alpha\in\{A,B\}}L^\alpha(\ubf^\alpha),
\end{equation}
and under these assumptions $E$ attains a minimizer on the admissible set $K$. The argument is the
direct method: polyconvexity gives weak lower semicontinuity of $E$, the growth bound keeps a
minimizing sequence bounded in $H^1\times H^1$, the trace inequality and the continuity of the gap
operator make $K$ weakly closed, and weak compactness then delivers the minimizer
\citep{ball1977convexity,ciarlet1988mathematical}. Adding the penalty changes nothing essential: the
penalty energy $\Pi_c$ of~\eqref{eq:penalty-potential} is non-negative and convex in $\gn$, so adding
it to $E$ preserves both lower semicontinuity and coercivity, and the penalized problem is well posed
for every $\varepsilon_n>0$.

% code: contact_atlas/03_mathematical_theory.md
The remaining question is how close the penalty solution lands to the exact one, and how that closeness
depends on the stiffness. Let $\ubf_\varepsilon$ be the penalty solution and $\ubf^\star$ the exact
(multiplier) solution. The standard estimate of \citet{kikuchi1988contact} bounds their difference in
terms of the penalty stiffness,
\begin{equation}\label{eq:penalty-estimate}
\norm{\ubf_{\varepsilon}-\ubf^\star}_{H^1}\ \le\ \frac{C}{\sqrt{\varepsilon_n}},
\qquad
\norm{t_{N,\varepsilon}-t_N^\star}_{H^{-1/2}(\Gamma_c)}\ \le\ C\sqrt{\varepsilon_n},
\end{equation}
so the displacement converges at rate $O(\varepsilon_n^{-1/2})$ while the pressure converges at
$O(\varepsilon_n^{1/2})$. These two rates pull in opposite directions, and balancing the penalty error
against the $O(h)$ discretization error fixes the stiffness at $\varepsilon_n\sim E/h$. The augmented
Lagrangian removes the trade-off altogether: the Uzawa iteration~\eqref{eq:al-pressure} contracts,
$\norm{\lambda^{k+1}-\lambda^\star}\le\rho\,\norm{\lambda^{k}-\lambda^\star}$ with $\rho<1$ for any
$\varepsilon_n>0$, so the multiplier converges to the exact pressure, $\lambda^k\to t_N^\star$, without
sending $\varepsilon_n\to\infty$ \citep{alart1991mixed,simo1992augmented}. Collecting the pieces, the
total error in the contact pressure splits into three terms,
$\norm{t_{N,h}-t_N^\star}\le C_1 h + C_2\varepsilon_n^{-1/2}+C_3\,\epsilon_{\mathrm{geom}}$, where
$\epsilon_{\mathrm{geom}}$ is the geometric approximation error of the oracle, the $L^\infty$ error of
the chart or of the level set. The first two terms are common to both detectors and shrink with mesh
and stiffness. The third term is where the atlas and the level set part company, and it is precisely
the term that Sections~\ref{sec:verif} and~\ref{sec:cv7} measure.
